% Vista editorial de corpus/source/masas/02_ley_general.tex, líneas 363--491: cuerpo exacto salvo el retítulo académico del epígrafe sobre el ciclo de 108 estados. Las líneas 1--263 ya residen exactamente en 72_rev2_nucleo_ley_i.tex y las líneas 264--362 en 72_rev2_algoritmo_universal.tex.
\chapter{Systematic extension of mass theory}
\Needspace{7\baselineskip}
\section{Unified foundation of mass}
\Needspace{7\baselineskip}
\subsection{The object prior to magnitude}
The theory begins before the number. A physically individuated mass is not a real number accompanied by the name of a particle, but the terminal term of a mathematical biography. APP supplies the finite incidence space and the two correlative operations; TRIT preserves the temporal orientation of each transition; TPK transports phase, sheet, carry, boundary, route, memory, and survival. Dimensional Mechanics subsequently converts that biography into a section of the mass line. The complete order is

\[
\mathrm{APP}\longrightarrow\mathrm{TRIT}\longrightarrow\mathrm{TPK}
\longrightarrow\Gamma^{\rm surv}
\longrightarrow\gamma_{\rm obs}
\longrightarrow\mathscr L_\gamma
\longrightarrow\sigma(\gamma)
\longrightarrow\chi_{\rm mass}
\longrightarrow\mathcal T_{90,120}
\longrightarrow\mathcal L_M.
\]
Each arrow retains information used by the next. The cell does not replace the route; the route does not replace the register; a signature does not replace its genealogical register; a character does not replace the dimensional section on which it acts. This separation explains why two particles may share charge or spin and nevertheless have different masses: their enriched biographies are not the same. It also explains why particle and antiparticle share mass: the involution changes odd orientation and preserves even memory, which descends to the mass quotient.

APP consists of nine by nine published positions. Each position admits four initial orientations, so the oriented domain contains

\[
81\cdot4=324
\]
candidate routes. The internal classification produces 56 route families and 13 multisections. These cardinalities are simultaneous: 81 counts cells, 324 counts orientations, 56 counts quotients by history, and 13 counts coarse operator types. No particle table replaces these domains.

\Needspace{7\baselineskip}
\subsection{Discrete unit of the continuum that supports the mass}
The mass line is not appended to five independent theories of the continuum. It follows from the unique APP--TRIT--TPK state from which the spectral, solenoidal, gauge, coherent, and determinantal aspects jointly emerge, designated in the construction notation by \(sp,sol,gau,coh,det\). Those symbols are reading discriminants, not prior subdomains. If \(\mathcal C_{\rm disc}\) denotes the common structure, the generating arrow has the form

\[
\mathfrak E_{\rm cont}:\mathcal C_{\rm disc}
\longrightarrow
\bigl(\mathcal C_\infty;
\mathcal O_{sp},\mathcal O_{sol},\mathcal O_{gau},
\mathcal O_{coh},\mathcal O_{det}\bigr),
\]
where the five observations belong to the same limit object, share fibers, orientation, carry, memory, boundary, and readers, and are transported by the same refinement squares. Mass uses its determinant reading for the action scale, its spectral reading for the character and the towers, its gauge reading for the null and charged sectors, and its coherent and solenoidal readings for propagation and composition. None of them is introduced by presupposing that the state belongs to a classical physical class.

This unit is a typing condition: separating one of those aspects and retaining the others would produce a different object, incapable of simultaneously supporting route, action, charge, amplitude, and mass. The present theory does not reprove the complete continuum; it incorporates its unified construction as a foundation and verifies that none of the mass insertions disaggregates it.

\Needspace{7\baselineskip}
\subsection{Absolute scale and quotients}
The local determinantal reader has Smith normal form.

\[
\operatorname{SNF}(H_4)=\operatorname{diag}(1,2,2,4),
\]
from which

\[
|\operatorname{coker}H_4|=16,
\qquad
\Sigma_H=\varphi_{\rm HMT}\alpha_{\rm HMT}^{16},
\qquad
\operatorname{ord}_{10}(\Sigma_H)=-34.
\]
The decade is obtained before the SI chart. The pre and return sections of the action line are

\[
\hbar_{\rm pre}=H_5\,10^{-34}u_S,
\qquad
\hbar_{\rm ret}=\eta_{\rm ret}^{(\deg)}\,10^{-34}u_S,
\]
and its bidirectional quotient is

\[
R_{\rm act}=\frac{H_5}{\eta_{\rm ret}^{(\deg)}}.
\]
In the clock branch, we explicitly choose the return section, \(\hbar_{\rm int}:=\hbar_{\rm ret}\), as the action section. A declared clock \(t_0\) provides

\[
\omega_0=\frac{2\pi}{108t_0},
\qquad
\varepsilon_{\rm clk}=\hbar_{\rm ret}\omega_0,
\qquad
m_{\rm clk}=\varepsilon_{\rm clk}c_{\rm int}^{-2}.
\]
Adopting \(u_M^{\rm clk}:=m_{\rm clk}\) as origin of the chart is a normalization explicit. Not identifies automatically this section with a basis energetic or mass previously declared. The character produces first the operator basal

\[
\widehat M_X^{(0)}
=\mathfrak e_0m_{\rm clk}\,
\chi_{\rm full}(\sigma_X-\sigma_e,\eta_X-\eta_e),
\]
and the bidirectional reader acts afterward. For each presentation \(r\in\{D,\Omega\}\), let \(\widehat B_X^r=(\widehat B_X^r)^*\) over the finite fiber; in an infinite realization one also requires a common dense domain for functional calculus. Then

\[
\widehat M_X^{\beta,r}
=R_{\rm act}^{\widehat B_X^r/2}
\widehat M_X^{(0)}
R_{\rm act}^{\widehat B_X^r/2}.
\]
On a scalar fibre, \(\widehat B_X^r=\kappa_XI\), so that

\[
m_X^\beta=m_X^{(0)}R_{\rm act}^{\kappa_X}.
\]
If the tower coordinates are already expressed relative to the electron, \(-\eta_e\) is omitted. When a quotient is formed, the common dimensional section cancels, but the difference in bidirectional reading does not:

\[
\frac{m_Y^\beta}{m_X^\beta}
=\chi_{\rm full}(\sigma_Y-\sigma_X,\eta_Y-\eta_X)
R_{\rm act}^{\kappa_Y-\kappa_X}.
\]
Therefore, \(10^{-34}\) fixes the absolute scale of action, energy, and mass; it is not a hidden relative corrector. The fine information of a species can enter only through its character, its relative clock, or the oriented quotient of the pre- and return sections.

\Needspace{7\baselineskip}
\subsection{Internal quantity of the mass line induced by the cycle of 108 states}
The Planck--CT108 self-duality provides, in the internal normalization

\[
\hbar_{\rm int}=c_{\rm int}=t_0=1,
\]
the section

\[
m_{108}^{\rm int}
=0.058177641733144319230789692282953757\ldots .
\]
This number is a normalized quantum of the internal mass line. It is not the mass of a particle, does not by itself carry a MeV or GeV chart, and cannot be inserted as an additional species row. Its function is to verify compatibility among the action chart, the CT108 clock, and the self-duality of the mass line. The transition to a nominal mass still requires the route, signature, operator, and dimensional section of the species.

\Needspace{7\baselineskip}
\subsection{Existence of the HMT--MD output for surviving extensions and scalar-reduction criterion by fiber rank}
Given a surviving extension, a cellular reader, an integer signature, an absolutely convergent family of tower coefficients, and a dimensional mass section, a well-typed HMT--MD output exists. If the fiber has a unique fixed point compatible with the involution, the output reduces to a canonical scalar. If the fiber has higher rank, the primary output is the spectral operator on the fiber; a physical scalar mass requires the coupling functional that individuates the observed representation.

This theorem does not weaken the recovered scalar evaluations. It distinguishes three claims: the internal operator is constructed; the evaluation of a declared signature is exact; and the prospective determination of that signature is an additional theorem when it is not deduced from the reader itself. Experimental comparison is performed only after fixing which of those three claims is being published.

\Needspace{7\baselineskip}
